\documentclass[11pt]{article}
\pdfoutput=1
\usepackage{amssymb,latexsym,amsmath,amsbsy}
\usepackage[dvips]{graphicx}
%\usepackage{Young}  % to draw young diagrams
\usepackage{cite}
%\usepackage[breaklinks=true,colorlinks=true,linkcolor=blue,urlcolor=blue,citecolor=blue]{hyperref}
\usepackage{hyperref}
%\usepackage[notref,notcite]{showkeys}
%\pagestyle{empty}
\headheight=0mm
\headsep=-20mm
\oddsidemargin=0mm
\evensidemargin=0mm
%\textheight=235mm
\textwidth=165mm
%
% definitions concerning automatic numbering of definitions, etc. %
\newtheorem{theorem}{Theorem}
\newtheorem{definition}[theorem]{Definition}
\newtheorem{lemma}[theorem]{Lemma}
\newtheorem{corollary}[theorem]{Corollary}
\newtheorem{remark}[theorem]{Remark}
\newtheorem{proposition}[theorem]{Proposition}
%
\DeclareMathOperator{\tr}{tr}
\DeclareMathOperator{\ch}{char}
\DeclareMathOperator{\sdim}{sdim}
%
\newcommand{\myatop}[2]{\genfrac{}{}{0pt}{}{#1}{#2}}
\newcommand{\BE}{\begin{equation}}
\newcommand{\EE}{\end{equation}}
\newcommand{\BQ}{\begin{equation} \begin{array}{c}}
\newcommand{\EQ}{\end{array}\end{equation}}
\newcommand{\BT}{\begin{theorem}}
\newcommand{\ET}{\end{theorem}}

\newcommand{\BL}{\begin{lstlisting}}
\newcommand{\EL}{\end{lstlisting}}


\usepackage{listings}
\usepackage{color}

\renewcommand\lstlistingname{Quelltext} % Change language of section name

\lstset{ % General setup for the package
	language=C,
	basicstyle=\small\sffamily,
	numbers=left,
 	numberstyle=\tiny,
	frame=tb,
	tabsize=4,
	columns=fixed,
	showstringspaces=false,
	showtabs=false,
	keepspaces,
	commentstyle=\color{red},
	keywordstyle=\color{blue}
}

\begin{document}
%\begin{center}
\noindent
{\Large \bf
The Magic2 deep RNA-sequenecing analyser\\User Guide
}\\
\\
{\bf G.~Boratyn}\\[2mm]
{\bf D.~Thierry-Mieg}\\[2mm]
{\bf J.~Thierry-Mieg}\\[2mm]
National Library of Medicine, National Institutes of Health,\\
8600 Rockville Pike, Bethesda MD20894, USA.\\
 E-mail: mieg@ncbi.nlm.nih.gov \\[2mm]
%\end{center}

%\vskip 10mm

%\noindent
% Short title: AQL

%\vskip 3cm
%\noindent
%Communicating author: J. Thierry-Mieg

%\noindent

%\noindent
%PACS numbers: 03.67.Hk, 02.30.Gp

\begin{abstract}
  Magic2 is an accurate and sensitive long and short reads RNA deep-sequencing analyser.
  It can read local fasta/fastq files or import data-sets on the fly from the NCBI short-read-archive (SRA).
  It can export in a signle pass the alignment files, the coverage plots, the gene expression and many metrics.
  It is very easy to use because it analyzes in a dry run the top of the sequence files and the hardware
  to select the optimal strategy. This guide explains how to use this program.

  
  Magic2 is freely avaliable from
  git@github.com:NLM-DIR/Magic2
\end{abstract}

%\vfill\eject
%\addtolength{\parskip}{2mm}

\newpage
\tableofcontents
\newpage

%%%%%%%%%%%%%%%%%%%%%%%%%%%%%%%%%%%%%%%%%%%%%%%%%%%%%%%%%%%%%%%%
% SECTION
%%%%%%%%%%%%%%%%%%%%%%%%%%%%%%%%%%%%%%%%%%%%%%%%%%%%%%%%%%%%%%%%
\section{Introduction}
Considering the fast evolution of the sequencing technologies, we need to
accuratelly aligned an ever greater number of short paired-end Illumina reads
and ever longer reads from PacBio, Nanopore, Roche and other manufacturers.
The files are very large and it is important to minimize the number of steps
and intermediate formats.

Magic2 can align local fasta/fastq files
but it can also automatically download
SRR runs from NCBI short read archives (SRA). 
Optionally, these data can be cached on the local computer as new fastq files, but
if you have a fast network connection to NCBI, this
is not mandatory. Converselly, Magic2 is convenient just to download the data
that you wish to analyze with other programs.

The program recognizes the eventual adaptors and reports very detailled quality control metrics,
It is fast because, in a single pass,
it can export the alignments, the gene expression counts,
the known and new introns and their support, 
the polyA addition sites.
and a collection of coverage plots indicating the gene expression,
the probable breakpoints and the candidate transcriptions starts and stops.

Section 1 gives simple examples of how to use Magic2.
Section 2 explains how to download and recompile the program.
Section 3 gives all the details on the recognized sequence formats and the way
to download the runs from SRA.
Section 4 explains the contruction of the target index.
The remaining section present the various output files and how to use them.

If available, the code will delegate parts of the calculations to a GPU.

\section {Simple examples}

First download the code on a Linux or Mac machine with at least 32 Giga bases or RAM.
The code will already run if you have 4 processors, and will run faster
on a larger machine because it automatically adapts its level of paralellism to the
host computer.

\subsection {Align data on a raw genome}

To just align a set of reads just on a genome, you need to type
two successive commands

\BL
magic2 --createIndex IDX.test  -t genome.fasta.gz
magic2 -x IDX.test -i read.fasta.gz -o out 
\end{lstlisting}

where 'genome.fasta.gz' is a single fasta file containing the sequence of all the chromosomes
and 'read.fasta.gz' contains the reads. The results will be visible in the
output directory 'out'. The data files can be gzip compressed or not, the sequences
can be given in fasta or fastq format.

To align paired end reads, the pairs can be given in two fasta/fastq files
separated by a plus sign, 
or in interleaved format
\BL
magic2 -x IDX.test -i read_R1.fastq.gz+read_R2.fastq.gz -o out2a
magic2 -x IDX.test -i read.sample_12.fastq.gz -o out2b
\end{lstlisting}

Finally, the data can be downloaded on the fly from the NCBI Short-Read-Archives (SRA)
and analyzed simply by providing a set of SRR identifiers

\BL
magic2 -x IDX.test -srr SRR1234,SRR1235,SRR1236 -o out23
\end{lstlisting}

Notice that you not need to say if the sequences are DNA or RNA, long or short, or if they
contain adaptors. All these properties are automatically recognized by the code by performing a dry
run on the first 10 Mega bases, the realigning everything with the propoer strategy.

\subsection {Align data on an annotated genome}

More accurate and more detailled result will be obtained if you provide a known annotation.
Magic2 will favaorize jumping the known introns, and will report the gene expression counts.
It will also report the fraction of the sequences mapping on the ribosomal RNA (sometimes as high as 80\%)
and on the mitochondria. This is important because a complete genome, like the human T2T project,
contains very many repetitions of the ribosomal RNAs and echos of parts of the mitochondria,
and many programs, including Magic2, would not report highly repetitive alignments.

To provide an annotation, you need to create a small target configuration file
called for example T2T.tConfig  containing 2 tab separated columns
\BL
G        T2T.chromosomes.fasta.gz
M        T2T.mitochondria.fasta.gz
R        T2T.rRNA.fasta.gz
I        T2T.gtf.gz
\end{lstlisting}
where T2T.gff.gz is the T2T annotation file, available for example from
NCBI or EBI in gtf or gff format.

Then, create an annotated index by using the parameter (-T) rather then (-t), and align
\BL
magic2 --createIndex IDX.T2T  -T  T2T.tConfig

magic2 -x IDX.T2T -srr SRR1234,SRR1235,SRR1236 --full -o out5
\end{lstlisting}
If you have a very large machine, you could process the 3 runs in parallel,
the memory needed by the index will be shared, but still you will consume larger ressources
\BL
magic2 -x IDX.T2T -srr SRR1234 --full -o out5a &
magic2 -x IDX.T2T -srr SRR1234 --full -o out5b &
magic2 -x IDX.T2T -srr SRR1236 --full -o out5c &
\end{lstlisting}
In the ouput directory, you will find the alignemnt files, the polyAs, the intron counts
and if you used the --full parameter, yo will also find the gene expression counts
the canditate transcription starts and ends and a collection of coverage plots.


\section {Downloading and installing Magic2}

\subsection {Simply download the executables}

The precompiled executable files can be found at
\BL
where are the binaries
\end{lstlisting}
select your machine (Mac or Linux) and eventually also download the GPU enabled code

\subsection{Download the source code and recompile}

The code is open-source and freely available from github.
\BL
git clone git@github.com:NLM-DIR/Magic2
\end{lstlisting}

In case of very intensive use,
the precompiled executables may be slightly slower than the programs recompiled locally
because the latter may be better taylored to the instruction sets of your actual hardware.
To recompile the C code. stay in the download directory and type
\BL
tcsh
cd ./Magic2
setenv ACEDB_MACHINE LINUX_4_OPT
make libs
cd wsa
make
\end{lstlisting}
The executable will be the file ./Magic2/bin.LINUX\_4\_OPT/magic2
Add this directory to your path or copy magic2 in /usr/local/bin
so that the command will be always found. You can test if the code run
by typing 'magic2 --help'
//
//
\section {Using sequence data available in SRA}
\subsection {Magic2 can download files from SRA}

Magic2 is a convenient tool to download runs from the NCBI short read archives
in a format recognized by most bioinformatics program.

Suppose that you are interested in the human paired-end Illumina run SRR24511885.
The command
\BL
magic2  --sraDownload SRR24511885     // creates SRA/SRR24511885.sample_12.fasta.gz
\end{lstlisting}
will download an interleaved fasta file, in which each pair occupies 4 successive lines
and the read identifiers will be s.number/1 and s.number/2.
If you are interested in the quality factors, add a parameter --fastq, and if you
want to split the pairs into 2 separate files use --split\_pairs.

The full syntax of the command is
\BL
magic2  --sraDownload SRR1,SRR2,SRR3,... [--fastq] [--split_pairs] --maxGB <int>
\end{lstlisting}
 and for example
\BL
magic2  --sraDownload SRR24511885 --fastq --split_pairs --maxGB 2
 // creates SRA/SRR24511885_R1.fastq.gz and SRA/SRR24511885_R2.fastq.gz  
\end{lstlisting}
downloads the first 2 Gigabases of run SRR24511885 and stores it as 2 fastq files.
By default, in the absence of the --maxGB parameter, the whole file is downloaded.
The number of SRR identifier in the list is not limited.
The code will automatically export a single file SRRn.fastq.gz for the single end runs.

\subsection {Magic2 can process on the fly SRA runs}

If you have a good connection to the NCBI servers, you do not need to store the files
locally on your computer. You can align them on the fly with the command
\BL
magic2  -x IDX --srr SRR1,SRR2,SRR3 -o out_directory 
\end{lstlisting}
magic2 will recognize single-end versus paired end runs and will start aligning
as soon as the first few magabases start flowing in.
Often, the download may be the limiting factor
and all the sequences will be aligned a few seconds after the end of the transfer.

An interesting way to use this command is to evaluate the quality cotrl metrics
on the first 2 Gigabases of a number of SRR runs without downloading the files and
without exporting the alignments
\BL
magic2  -x IDX --srr SRR1,SRR2,SRR3 --no_ali -o out_directory 
\end{lstlisting}
This will produce quite fast a detailled portrait of all these runs without
consuming any appreciable disk space on your computer.

If the files are already present in the directory ./SRA, for example if they
were explicitly downloaded by the previous commands, these cached files will
be used rather than accessing the network.

\subsection {Magic2 can align and cache the data at the same time}

You can combine the previous commands by giving at the same time a list of SRR identifiers
a caching request and a target index. For example
\BL
magic2  --srr SRR24511885 --fastq --sraCaching -x IDX.T2T -o out_dir --full
// creates SRA/SRR24511885.sample_12.fastq.gz
// produces the complete analysis of the run in directory out_dir
\end{lstlisting}
If the SRR run is already in the ./SRA cache directory  in any of the recognized formats,
it will not be downloaded a second time.

\section {Creation of a target index}

As explained in the introduction, a given run can be aligned on a raw genome,
but it is often more interesting to compared the reads to an annotated genome.
The first step is to construct a taget configuration file

\subsection {The target configuration file tConfig}

Using a text editor, construct a configuration file describing your target.
The format is simple. There are 2 tab separated columns. The first
column contains a single letter giving the type, the second column contains a file name.

The recognized types are G, M, C, R, E, A, I, B, V. Several lines can start with the same letter.

G: genome, the fasta file give the genomic sequence of the target,
but if more convenient, you can provide each chromosome fasta file on a separate line starting with G

M: Mitochondria, the fasta file of the mitochondria

C: In plants, the sequence of the chloroplast

R: the sequence of the ribosomal RNA precursors. Giving these sequences this is very important because
ribosomal RNA may sometimes represent a very high fraction (80%) of a sample
and the sequence is repeated many times in the chromosome creating confusion


E: ERCC and other control sequences. Also give the DNA control

A: Adaptors fasta sequence, this file is optional, adaptors can be recognized on the fly

B:  bacteria, possible contamination 

V: virus, possible contamination 

I: Annotation file in gtf or gff format, or a list of introns

In each case the the chromosome name (column 1) must match the G fasta file(s).

\BL
In gtf/gff format the code expects (at least) 7  columns chrom/method/tag/a1/a2/./[+|-]\n"
tag : [exon] other lines are ignored
a1 a2 : the positions of the first and last base of the exon.
strand : [+|-] indicates the strand, a1 < a2, and the coordinates are 1-based.

In .introns format the code expects 3 columns chrom/a1/a2
a1 a2 : give the positions of the 2 G of the Gt_aG motif
on the top strand a1 < a2, on the bottom strand a1 > a2.
\end{lstlisting}

Please create the .introns file if the gff/gtf file is not available

\subsection {Creating the index}

Magic2 belongs to the seed/extend class of aligners.
In a first pass, magic2 extract from the genome a set of seeds of fixed length
and then searches these seeds in the reads. The length and the density of the genome seeds
may be adjusted. By default, magic2 use 16-mers up to 200 Megabases (drosophila)
and 18-mers for longer genomes (human). The default stepping is s=6, meaning we choose the
'best' seed among the next s words. As a result, the maximal distance between seeds is s,
and the average distance is s/2, exactly as when using the special seeds refered by Durbin,
but in a much simpler way.

The binary files containing the collected seeds are stored in the index directory, together with
a summary of the annotaton files. To contruct the index please run the command

\BL
magic2 --createIndex IDX.target -T tConfig.target [--seedLength <int>] [--step <int>]
\end{lstlisting}
Changing the default seed length is not recommanded in metazoaires, 
but it may possibly be interesting
to choose shorter seed in a bacteria project.
Using a longer stepping like 8 or 10  will produce a faster code, consuming less RAM,
but will reduce the sensitivity of the aligner.



\section {Alignments}

In this project, our main goal is completeness and accuracy. We hope to align
the highest possible percentage of reads and of bases in each read, even in
the presence of sequencing noise.

To align a run, provide an SRR identifier (it will be downloaded automatically
as explained above), or give a fasta or a fastq file. For paired end runs,
if they are downloaded from SRA, they are recogbized automatically,
other wise you can provide either a pairs of fasta/fastq files separated by a plus sign or
provide an interleaved fasta/fastq file which must be called xxx.sample\_12.fasta[aq][.gz]

\BL
magic2 -x IDX.target -i f1.fasta.gz+f2.fasta.gz -o out_dir
\end{lstlisting}

By default, the alignments are exported in the magic .hits format, but they can also be 
exported in the magic-blast tabular format or in bam format by providing the parameters

\BL
magic2 -x IDX.target -i f1.fasta.gz+f2.fasta.gz -o out_dir --hits --tabular --bam
\end{lstlisting}
If several formats are requested, several files will be exported. The bam format is well known,
the .hits and .tabular format give more details, in particular they discribe the
mismatches and the introns, and the overhangs more explicitely.

One may also block the exportation of the alignments with the parameter --no\_ali
\BL
magic2 -x IDX.target -i f1.fasta.gz+f2.fasta.gz -o out_dir --no_ali
\end{lstlisting}

\section {Intron support}

Our secondary goal is to simplify the user's experience. The program
does not require any cost or strategy parameters. It evaluate by itself, via
a dry run on the first 10 mega bases, the sequencing technology and the
characteristic of the sample and automatically recognizes and clips
the eventual adapters. It also evaluates the hardware to decide the optimal level
of parallelisation and if available will run part of the calculations on a GPU.

Thanks to its non-standard parallelization by GO like channels,
computing the coverage plots and the gene expression counts come at essentlially no
cost, whereas in current  methods (HISAT, Star etc.)  they usually require a very slow sorting and rescanning
of the BAM files. Actually, in Magic2, exporting the BAM alignment files can be considerd
optional, since Magic2 may be abale to export directly all the data useful to your project.

A collateral benefit of Magic2 is its speed. The computers, over the last ten to twenty years,
have considerably accelerated the computing inits (CPUs), but memory access kept lagging.
To compensate, multi-layered cache architectures were designed. The CPUs are grouped in cores,
and access the large mamory (the RAM in giga bytes) via several level of caches (typically L1 in kb, L2 in Mb).
Considering that memory access is the bottle neck, we designed the Magic2 algotrithms to priviledge
sequential memory access. For example

\section{Conclusion}


%%%%%%%%%%%%%%%%%%%%%%%%%%%%%%%%%%%%%%%%%%%%%%%%%%%%%%%%%%%%%%%%
% SECTION
%%%%%%%%%%%%%%%%%%%%%%%%%%%%%%%%%%%%%%%%%%%%%%%%%%%%%%%%%%%%%%%%

\section*{Acknowledgments}

We would like to thank David Landsman and Tom Madden.
This research was supported by the Intramural Research Program of the National Library of Medicine, National Institute of Heath.

%%%%%%%%%%%%%%%%%%%%%%%%%%%%%%%%%%%%%%%%%%%%%%%%%%%%%%%%%%%%%%%%
% SECTION
%%%%%%%%%%%%%%%%%%%%%%%%%%%%%%%%%%%%%%%%%%%%%%%%%%%%%%%%%%%%%%%%
\begin{thebibliography}{99}

\bibitem {DTM92}
Richard Durbin and Jean Thierry-Mieg
Syntactic Definitions for the ACEDB Database Manager
Technical Report: MRC Lab. for Molecular Biology, 1992

\bibitem {DTM94}
Richard Durbin and Jean Thierry-Mieg
The acedb genome database,
Computational Methods In Genome Research, pages 45-55.
Edited by S. Suhai, Plenum Press, New York, 1994

\bibitem {STM98}
Lincoln D Stein, Jean Thierry-Mieg
Scriptable access to the Caenorhabditis elegans genome sequence and other ACEDB databases
Genome research, 8,12 pp 1308-1315, 1998

\bibitem {JDTMS}
Jean Thierry-Mieg, Danielle Thierry-Mieg and Lincoln Stein
ACEDB: The Ace Database Manager
In Databases and Systems, 2001
DOI: 10.1007/0-306-46903-0\_23


\bibitem{cDNA}
Danielle Thierry-Mieg and Jean Thierry-Mieg
AceView: a comprehensive cDNA-supported gene and transcripts annotation.
Genome Biology 2006, 7(Suppl 1):S12.  

\bibitem{Code}
Code and documentation are available from  
git clone git@github.com:NLM-DIR/Magic2

\end{thebibliography}


%%%%%%%%%%%%%%%%%%%%%%%%%%%%%%%%%%%%%%%%%%%%%%%%%%%%%%%%%%%%%%%%
% SECTION
%%%%%%%%%%%%%%%%%%%%%%%%%%%%%%%%%%%%%%%%%%%%%%%%%%%%%%%%%%%%%%%%
\appendix
\section {Annex}
%%%%%%%%%%%%%%%%%%%%%%%%%%%%%%%%%%%%%%%%%%%%%%%%%%%%%%%%%%%%%%%%
% SECTION
%%%%%%%%%%%%%%%%%%%%%%%%%%%%%%%%%%%%%%%%%%%%%%%%%%%%%%%%%%%%%%%%
\subsection{List of operators and reserved keywords}

To summarize, the full list of operators, in their order of precedence is
\BL
   ; TITLE order\_by
   from select , where
   in := = 
   || OR ^^ XOR &&  AND ! NOT
   like =~ ~ == != >= <= > < >~ <~
   ISA
   modulo + - * / ^
   DNA PEPTIDE PEP CODING DATEDIFF 
   COUNT MIN MAX SUM AVERAGE STDDEV
   .class .name .timestamp
   >> -> # => 
   @ OBJECT
   () {} []
\end{lstlisting}

with the restriction that the 3 kinds of parentheses and 
also the quotes must be nested correctly  (["bb"]) is correct but (x[y)] is illegal. Operators written in capitals, i.e. DNA or AVERAGE, must be written in capitals (are case-sensitive)

%%%%%%%%%%%%%%%%%%%%%%%%%%%%%%%%%%%%%%%%%%%%%%%%%%%%%%%%%%%%%%%%
% SECTION
%%%%%%%%%%%%%%%%%%%%%%%%%%%%%%%%%%%%%%%%%%%%%%%%%%%%%%%%%%%%%%%%
\subsection{History of the development of the acedb query language}

The query syntax described in this document results from the convergence of several developments
   - the original acedb query language of 1990, available in the 'query' box of the acedb graphic interface
   - the table maker system of 1992, with its user friendly graphic interface 
   - the original AQL query language developed at the Sanger around 2002

The general idea was to learn from the 3 existing systems, conserve the best aspects of each, fix their limitations and clean up the interface. For example, 

-- The acedb query language is very terse and expressive, but it is limited to constructing sets of objects and does not allow to display the content of selected fields.

-- The table maker was meant to fix these limitations. Using its rich graphic interface, one can easily construct a tables involving objects, numbers, texts and even DNA sequences. However, even a simple tables cannot be constructed on the command line. One need always to construct the table graphically, to save its definition in a definition file, and run the query using that file. 

-- The original AQL query language was ment to fix this limitation. AQL defined a command line syntax, reminiscent of SQL, but adapted to the idiosyncrasies of an object oriented database like acedb. It made it in principle possible to compose a table on the command line or using a library call. Unfortunately, the syntax was slightly too rich, encouraging the user to write convoluted queries, and always verbose. Since there was no graphic interface to help compose a valid query, and since the original AQL compiler did not explain the eventual errors, it was hard to write a non-trivial valid query. Finally, AQL execution was slower than table-maker, and lacked access to DNA and protein sequences. 

From the user point of view, the previous interfaces are maintained, but the old style queries are translated internally and executed by the new query engine. In particular, the graphic table-maker interface remains valid and can still be used to construct complex queries and most original AQL queries remain valid. In practice the shortcuts 
defined in section 2.5 map the terse query language of the acedb graphic interface, first released in 1990, into a (modification of) the full fledged 
AQL queries developed around 2002. As a result, these 2 types of queries, and the table constructed via the acedb graphic table 
maker interface, are all expressed in the new unified syntax and the database exploration can be executed using a single highly optimized query engine.


%%%%%%%%%%%%%%%%%%%%%%%%%%%%%%%%%%%%%%%%%%%%%%%%%%%%%%%%%%%%%%%%
% SECTION
%%%%%%%%%%%%%%%%%%%%%%%%%%%%%%%%%%%%%%%%%%%%%%%%%%%%%%%%%%%%%%%%
\subsection {Code availability}

The acedb source code can be downloaded from 
 
$https://github.com/ncbi/AceView$, $https://github.com/jtmieg/AceView$ 

The tar ball is effectively identical to the distribution of the AceView/Magic pipeline 
\cite{DJTM} which uses acedb as the underlying database, and can organize the workflow and manage the
meta-data of tens of thousands of next-generation sequencing runs. The same package
underlies the Gene annotation public web site https://www.ncbi.nlm.nih.gov/AceView \cite{cDNA}

The program is written in C and has very few external dependencies. To compile the code try
\BL
  tcsh 
  gunzip -c magic.*.tar.gz | tar xf -
  setenv ACEDB_MACHINE LINUX_4_OPT
  cd magic*
  make libs tace
  make -k all
\end{lstlisting}
 
The code can also compile on the Mac using 
setenv ACEDB\_MACHINE  MAC\_X\_64\_OPT
or on any other linux/unix platform. Multiple choices are offered in the directory wmake.
The tace command-line interface presented in the present document
should compile without difficulty.
The xace graphic-interface requires 
the installation of the public X11 development environment
as explained in the README file situated at the top of the acedb distribution.

A test-bench containing examples of queries can be constructed
using the command
\BL
  tcsh wbql/bqltest.tcsh
\end{lstlisting}
After running the script, 
the file ACEDB\_QUERY\_LANGUAGE\_TEST/test.out contains the output of
the commands contained in test.cmd and the file test.diff,
which contains the differences between test.out and test.out.expected, should be empty.
The database can be examined and
further queries can be tested using the command
\BL
  bin*/tace ACEDB_QUERY_LANGUAGE_TEST
\end{lstlisting}
which will activate the command line interface of acedb. All the examples
in the present document can be copied in the interface and
should give proper answers. A question mark will invoke
the on-line help and list all the possible commands.
In particuler, the acedb command 'show' will display
the full content of the active objects and will help to understand why 
they where selected by the query. 

Please try the code and send us feedback. 



\subsection {Running an acedb session}

As explained above, a small test database can be constructed
automatically.
More generally, the directory waligner/metadata contains 
some examples of large acedb schemas.
Once you have constructed a local directory
containing an empty sub-directory database and a sub-directory wspec
defining the schema and a data file xxx.ace, you can
initialize a database and parse an ace file as follows:
\BL
  bin*/tace .  // launch the compiled text-acebd on the local directory
y              // yes, initialize the database
   parse myfile.ace              // parse some data
   save  // save the data in compiled binary format 
   quit
\end{lstlisting}
 Acedb can be regarded as a data compiler. It can read successively several
data files, merge them and compile them into a compact binary format allowing
fast retrieval. Clear messages signal places in the data file which do not conform
to the schema. You can fix the file and reload it. Reading the same file several
times is idempotent, i.e. it does not cause any problem.

Subclasses are configured in the file wspec/subclasses.wm

\BL
   Class Prolific_author      // The class name must start with a letter
   Visible                    // Make it available in graphic menus
   Is_a_subclass_of Person    // Inherit from class Person
   Filter "COUNT Papers >= 5" // Chosen filter
\end{lstlisting}

One can then use the database and ask queries on the command
line or import it from a file using the $-$$i$ parameter
or redirect the output to a file using $-$$o$
\BL
  bin*/tace . // start the system
    select ...  from ... where // a query
xx yy zz      // the reply comes on stdout
xx yy zz
    select -o f.out ... from ... where ...
              // the reply goes in the file f.out
    select -a ... // the reply is quote protected (.ace format)
    select -i f.bql ... // file f.bql contains the query
    select -i f.bql -s  // silent mode
    select @ ... // use the new active set 
    ? select  // help of the command line parameters
    ?         // list all the acedb commands
    quit
\end{lstlisting}
In silent mode (with semi column at the end of the query, or using
the option -s (-silent)) the output is suppressed but the
active list is modified.

Acedb contains dozens of commands available by typing ? or ?command.
There is also a rich graphic interface, available using the
command $bin.*/xace .$, which generates on the fly in HTML5/SVG language
all the plots and diagrams presented on the AceView server \cite{cDNA}.
It provides among other tools
a graphic interface helping to compose complex tables.
But the description of the graphic acedb system is out of the scope of this 
user guide.

The MAGIC pipeline  \cite{DJTM} is a good example of a complex system 
relying on acedb.



%%%%%%%%%%%%%%%%%%%%%%%%%%%%%%%%%%%%%%%%%%%%%%%%%%%%%%%%%%%%%%%%
% SECTION
%%%%%%%%%%%%%%%%%%%%%%%%%%%%%%%%%%%%%%%%%%%%%%%%%%%%%%%%%%%%%%%%
\subsection{C language API}

The acedb system has a C language API. The application code
can either be linked into the open-source database server code,
in which case it directly shares the caches of the database engine,
or it can be part of a client code which calls the server
via tcp. The interesting point is that exactly the same code
works in both situations, one changes from a server to a client
code by simply changing the library used at link stage.
Here is an example of a complete C program invoking the query language  

\BL
#include <ac.h>  /* acedb C language API header */

int main (int argc, const char **argv)
{
  AC_HANDLE h = ac_new_handle () ;
  const char *errors = 0 ;
  const char *dbNameS = "acedb_directory" ;              // server case
  const char *dbNameC = "t:machine_name:port_number" ;   // client case
  AC_DB db = ac_open_db (dbNameX, &errors);              // X = S|C
  const char *query = "select p in class person where p like a*" ;
  AC_KEYSET aks = 0 ;  /* initial keyset, populates the @ symbol */
  AC_TABLE t = ac_bql_table (db, query, aks, errors, h) ;

  printf ("Found % lines in the table\n", t->rows) ;

  ac_db_close (db) ;
  ac_free (h) ;
  return 0 ;
} /* main */

\end{lstlisting}
The interface is detailled in the self documented header file wh/ac.h.

All comments are welcome.

%%%%%%%%%%%%%%%%%%%%%%%%%%%%%%%%%%%%%%%%%%%%%%%%%%%%%%%%%%%%%%%%
% SECTION
%%%%%%%%%%%%%%%%%%%%%%%%%%%%%%%%%%%%%%%%%%%%%%%%%%%%%%%%%%%%%%%%


%%%%%%%%%%%%%%%%%%%%%%%%%%%%%%%%%%%%%%%%%%%%%%%%%%%%%%%%%%%%%%%%

\end{document}
